% GENERATED FILE. Edit canonical lessons, then run npm run generate:books.
% canonical-source-hash: fnv1a64:d6931092b21a7127
% canonical-lessons: SA-C364-visesajnah, SA-C364-vimarsah, SA-C364-tika, SA-C364-nibandhah, SA-C364-nirupayati

\chapter{Ask an Expert}
\label{ch:sa-ask-an-expert}
\hlchaptermodality{364}

\begin{chapteropening}
\textbf{By the end of this chapter:} \emph{I can say who knows a subject best and what they found.}

Complete the last lesson of chapter 364: I can say who knows a subject best and what they found.
\end{chapteropening}

\section[\texorpdfstring{\sk{विशेषज्ञः} (viśeṣajñaḥ)}{विशेषज्ञः (viśeṣajñaḥ)}]{\sk{विशेषज्ञः} (viśeṣajñaḥ) — an expert}

\label{lesson:SA-C364-visesajnah}



\begin{quote}
Before the new one: say the Sanskrit for studies, then the Sanskrit for practises.
\end{quote}

\subsection*{You'll want to know: \sk{विशेषज्ञः}}
\textbf{\sk{विशेषज्ञः}} — \emph{viśeṣajñaḥ} — \textquotedblleft{}an expert\textquotedblright{}.

\begin{quote}
\sk{रोगस्य} \sk{विषये} \sk{विशेषज्ञं} \sk{पृच्छतु।}
\end{quote}

\begin{quote}
\emph{rogasya viṣaye viśeṣajñaṃ pṛcchatu.}
\end{quote}

\begin{quote}
Ask an expert about the illness.
\end{quote}

A masculine noun.

\begin{cousinweb}
\textbf{\sk{विशेषज्ञः}} (\emph{viśeṣajñaḥ}) is built from \emph{viśeṣa}, \textquotedblleft{}a distinction\textquotedblright{}, and \emph{jña}, \textquotedblleft{}knowing\textquotedblright{}; MW gives \emph{viśeṣa-jña} \textquotedblleft{}knowing distinctions, judicious\textquotedblright{}, and \textquotedblleft{}an expert\textquotedblright{} is its modern sense (MW).
\end{cousinweb}

\begin{grammarlens}[title={What you've built}]
One more word for giving reasons and opinions.
\end{grammarlens}

\subsection*{Your turn}
\begin{itemize}
\raggedright
  \item \emph{Say it:} \emph{viśeṣajñaḥ}
  \item \emph{Say it:} \emph{viśeṣajñaḥ}, once more
  \item \emph{Say it:} \emph{rogasya viṣaye viśeṣajñaṃ pṛcchatu.} — Ask an expert about the illness.
  \item \emph{Recall:} say the Sanskrit for studies, then the Sanskrit for practises, then say \emph{viśeṣajñaḥ} again
\end{itemize}

\begin{tcolorbox}[breakable,colback=teal!4,colframe=teal!35!black,title={Before you move on}]
What does \sk{विशेषज्ञः} mean? (\textquotedblleft{}An expert\textquotedblright{}.) Say it once more.
\end{tcolorbox}
\section[\texorpdfstring{\sk{विमर्शः} (vimarśaḥ)}{विमर्शः (vimarśaḥ)}]{\sk{विमर्शः} (vimarśaḥ) — deliberation}

\label{lesson:SA-C364-vimarsah}



\begin{quote}
Before the new one: say the Sanskrit for practises, then the Sanskrit for an expert.
\end{quote}

\subsection*{You'll want to know: \sk{विमर्शः}}
\textbf{\sk{विमर्शः}} — \emph{vimarśaḥ} — \textquotedblleft{}deliberation\textquotedblright{}.

\begin{quote}
\sk{विमर्शं} \sk{विना} \sk{निर्णयं} \sk{न} \sk{करोतु।}
\end{quote}

\begin{quote}
\emph{vimarśaṃ vinā nirṇayaṃ na karotu.}
\end{quote}

\begin{quote}
Make no decision without deliberation.
\end{quote}

A masculine noun.

\begin{cousinweb}
\textbf{\sk{विमर्शः}} (\emph{vimarśaḥ}) is built on \emph{vi} and the root √mṛś, \textquotedblleft{}to touch, to consider\textquotedblright{}; MW gives \emph{vi-marśa} \textquotedblleft{}consideration, deliberation, critical test, examination\textquotedblright{} (MW).
\end{cousinweb}

\begin{grammarlens}[title={What you've built}]
One more word for giving reasons and opinions.
\end{grammarlens}

\subsection*{Your turn}
\begin{itemize}
\raggedright
  \item \emph{Say it:} \emph{vimarśaḥ}
  \item \emph{Say it:} \emph{vimarśaḥ}, once more
  \item \emph{Say it:} \emph{vimarśaṃ vinā nirṇayaṃ na karotu.} — Make no decision without deliberation.
  \item \emph{Recall:} say the Sanskrit for practises, then the Sanskrit for an expert, then say \emph{vimarśaḥ} again
\end{itemize}

\begin{tcolorbox}[breakable,colback=teal!4,colframe=teal!35!black,title={Before you move on}]
What does \sk{विमर्शः} mean? (\textquotedblleft{}Deliberation\textquotedblright{}.) Say it once more.
\end{tcolorbox}
\section[\texorpdfstring{\sk{टीका} (ṭīkā)}{टीका (ṭīkā)}]{\sk{टीका} (ṭīkā) — a commentary}

\label{lesson:SA-C364-tika}



\begin{quote}
Before the new one: say the Sanskrit for an expert, then the Sanskrit for deliberation.
\end{quote}

\subsection*{You'll want to know: \sk{टीका}}
\textbf{\sk{टीका}} — \emph{ṭīkā} — \textquotedblleft{}a commentary\textquotedblright{}.

\begin{quote}
\sk{पण्डितः} \sk{ग्रन्थस्य} \sk{टीकां} \sk{लिखति।}
\end{quote}

\begin{quote}
\emph{paṇḍitaḥ granthasya ṭīkāṃ likhati.}
\end{quote}

\begin{quote}
The scholar writes a commentary on the book.
\end{quote}

A feminine noun.

\begin{cousinweb}
\textbf{\sk{टीका}} (\emph{ṭīkā}) MW gives \emph{ṭīkā} \textquotedblleft{}a commentary\textquotedblright{}, especially one on another commentary, and gives no derivation, so its origin is uncertain (MW).
\end{cousinweb}

\begin{grammarlens}[title={What you've built}]
One more word for giving reasons and opinions.
\end{grammarlens}

\subsection*{Your turn}
\begin{itemize}
\raggedright
  \item \emph{Say it:} \emph{ṭīkā}
  \item \emph{Say it:} \emph{ṭīkā}, once more
  \item \emph{Say it:} \emph{paṇḍitaḥ granthasya ṭīkāṃ likhati.} — The scholar writes a commentary on the book.
  \item \emph{Recall:} say the Sanskrit for an expert, then the Sanskrit for deliberation, then say \emph{ṭīkā} again
\end{itemize}

\begin{tcolorbox}[breakable,colback=teal!4,colframe=teal!35!black,title={Before you move on}]
What does \sk{टीका} mean? (\textquotedblleft{}A commentary\textquotedblright{}.) Say it once more.
\end{tcolorbox}
\section[\texorpdfstring{\sk{निबन्धः} (nibandhaḥ)}{निबन्धः (nibandhaḥ)}]{\sk{निबन्धः} (nibandhaḥ) — an essay; a composition}

\label{lesson:SA-C364-nibandhah}



\begin{quote}
Before the new one: say the Sanskrit for deliberation, then the Sanskrit for a commentary.
\end{quote}

\subsection*{You'll want to know: \sk{निबन्धः}}
\textbf{\sk{निबन्धः}} — \emph{nibandhaḥ} — \textquotedblleft{}an essay; a composition\textquotedblright{}.

\begin{quote}
\sk{छात्रः} \sk{धर्मस्य} \sk{विषये} \sk{निबन्धं} \sk{लिखति।}
\end{quote}

\begin{quote}
\emph{chātraḥ dharmasya viṣaye nibandhaṃ likhati.}
\end{quote}

\begin{quote}
The student writes an essay about dharma.
\end{quote}

A masculine noun.

\begin{cousinweb}
\textbf{\sk{निबन्धः}} (\emph{nibandhaḥ}) is built on \emph{ni}, \textquotedblleft{}down\textquotedblright{}, and the root √bandh, \textquotedblleft{}to bind\textquotedblright{}; MW gives \emph{ni-bandha} \textquotedblleft{}binding\textquotedblright{} and \textquotedblleft{}any literary composition or work\textquotedblright{}, and \textquotedblleft{}an essay\textquotedblright{} is the sense modern Sanskrit gives it (MW).
\end{cousinweb}

\begin{grammarlens}[title={What you've built}]
One more word for giving reasons and opinions.
\end{grammarlens}

\subsection*{Your turn}
\begin{itemize}
\raggedright
  \item \emph{Say it:} \emph{nibandhaḥ}
  \item \emph{Say it:} \emph{nibandhaḥ}, once more
  \item \emph{Say it:} \emph{chātraḥ dharmasya viṣaye nibandhaṃ likhati.} — The student writes an essay about dharma.
  \item \emph{Recall:} say the Sanskrit for deliberation, then the Sanskrit for a commentary, then say \emph{nibandhaḥ} again
\end{itemize}

\begin{tcolorbox}[breakable,colback=teal!4,colframe=teal!35!black,title={Before you move on}]
What does \sk{निबन्धः} mean? (\textquotedblleft{}An essay; a composition\textquotedblright{}.) Say it once more.
\end{tcolorbox}
\section[\texorpdfstring{\sk{निरूपयति} (nirūpayati)}{निरूपयति (nirūpayati)}]{\sk{निरूपयति} (nirūpayati) — investigates; determines}

\label{lesson:SA-C364-nirupayati}



\begin{quote}
Before the new one: say the Sanskrit for a commentary, then the Sanskrit for an essay.
\end{quote}

\subsection*{You'll want to know: \sk{निरूपयति}}
\textbf{\sk{निरूपयति}} — \emph{nirūpayati} — \textquotedblleft{}investigates; determines\textquotedblright{}.

\begin{quote}
\sk{विशेषज्ञः} \sk{रोगस्य} \sk{कारणं} \sk{निरूपयति।}
\end{quote}

\begin{quote}
\emph{viśeṣajñaḥ rogasya kāraṇaṃ nirūpayati.}
\end{quote}

\begin{quote}
The expert determines the cause of the illness.
\end{quote}

Past (he …): \sk{सः} \sk{कारणं} \sk{निरूपितवान्}, \emph{saḥ kāraṇaṃ nirūpitavān}.

\begin{cousinweb}
\textbf{\sk{निरूपयति}} (\emph{nirūpayati}) is built on \emph{ni} and \emph{rūpa}, \textquotedblleft{}form\textquotedblright{}, as a verb; MW gives \emph{ni-rūpita} \textquotedblleft{}observed, considered, weighed, ascertained, determined\textquotedblright{} (MW).
\end{cousinweb}

\begin{grammarlens}[title={What you've built}]
One more word for giving reasons and opinions.
\end{grammarlens}

\subsection*{Your turn}
\begin{itemize}
\raggedright
  \item \emph{Say it:} \emph{nirūpayati}
  \item \emph{Say it:} \emph{nirūpayati}, once more
  \item \emph{Say it:} \emph{viśeṣajñaḥ rogasya kāraṇaṃ nirūpayati.} — The expert determines the cause of the illness.
  \item \emph{Recall:} say the Sanskrit for a commentary, then the Sanskrit for an essay, then say \emph{nirūpayati} again
\end{itemize}

\begin{tcolorbox}[breakable,colback=teal!4,colframe=teal!35!black,title={Before you move on}]
What does \sk{निरूपयति} mean? (\textquotedblleft{}Investigates; determines\textquotedblright{}.) Say it once more.
\end{tcolorbox}
